% Generated draft: requires mathematical and source review.
\documentclass[11pt]{article}
\usepackage[margin=25mm]{geometry}
\usepackage{fontspec}
\setmainfont{Latin Modern Roman}
\usepackage{amsmath,amssymb,mathtools,unicode-math}
\setmathfont{Latin Modern Math}
\usepackage{enumitem,needspace,xurl,hyperref}
\setlength{\parindent}{0pt}
\setlength{\parskip}{5pt}
\setlength{\emergencystretch}{3em}
\newcommand{\N}{\mathbb{N}}
\newcommand{\Z}{\mathbb{Z}}
\newcommand{\Q}{\mathbb{Q}}
\newcommand{\R}{\mathbb{R}}
\newcommand{\C}{\mathbb{C}}
\newcommand{\F}{\mathbb{F}}
\newcommand{\A}{\mathbb{A}}
\newcommand{\PP}{\mathbb{P}}
\newcommand{\E}{\mathbb{E}}
\newcommand{\eps}{\varepsilon}
\newcommand{\dd}{\,\mathrm{d}}
\begin{document}
\section*{UnsolvedMath ID: 11000324}
\textbf{Source problem code:} catalogue entry
\subsection{Definitions and mathematical statement}
Let \(\mathcal H_{g,1}\), the maps \(\sigma,\overline\sigma\) and the degree-two classes \(\widetilde t_{2k+1}\) be as defined in the preceding source discussion. The notation here refers to the homology-cylinder/homology-cobordism setting indicated by that discussion, not the handlebody subgroup denoted similarly elsewhere. The precise quotient/group, map targets, class definitions and coefficients are absent.

The conjecture has two independent assertions, for every admissible \(k\):
\[\overline\sigma^{\,*}(\widetilde t_{2k+1})\ne0
\quad\text{in }H^2(\mathcal H_{g,1}),\]
\[\sigma^*(\widetilde t_{2k+1})=0
\quad\text{in }H^2(\mathcal H_{g,1}).\]
The bar distinguishes different maps and must not be dropped. The two displayed pullbacks do not assert that the same map sends a class both to zero and to a nonzero class. The parameter range, any required genus bounds and normalization of the trace classes must be supplied; ``any \(k\)'' is not silently chosen as \(k\geq0\) rather than \(k\geq1\).

In a usual homology-cylinder model, the objects are suitably boundary-marked three-dimensional cobordisms between copies of the genus-\(g\) one-boundary surface, with the inclusions inducing homology isomorphisms, and stacking gives the operation; the group \(\mathcal H_{g,1}\) typically uses a further homology-cobordism quotient. This only explains the contextual category, not the source's exact equivalence or the maps \(\sigma,\overline\sigma\).

The source calls the first assertion the group version of Conjecture 6.1 and expects it to be harder; that conjecture is absent. It says that nonzero \(\sigma^*\)-classes would give invariants of certain four-manifolds arising as two-dimensional families of homology cylinders, which the author views as unlikely and uses to motivate the second assertion. This is motivation, not a theorem prohibiting four-manifold invariants. The following announced problem is not supplied. The summary relates the discussion to Morita trace classes and the first MMM class/Casson secondary-invariant context, without an exact definition or argument settling either pullback.
\subsection{Short English statement}
For each allowed odd trace class, the barred pullback should be nonzero and the unbarred pullback zero in the homology-cylinder group's degree-two cohomology. The two maps, coefficients and parameter ranges are missing and must remain distinct.
\subsection{Sources}
\begin{itemize}
\item Local UnsolvedMath registry, permanent ID \(11000324\), exact source code \texttt{catalogue entry}.
\item Farb, \emph{Problems on Mapping Class Groups and Related Topics}, Conjecture 11.3, PDF page 379, with reference to Conjecture 6.1 and a following problem: \url{https://www.math.uchicago.edu/~farb/papers/mcgbook.pdf}. Registry access 2026-07-29; pdf-font-structure extraction.
\item The research summary mentions Morita's trace/MMM/Casson work, Sakasai's homology-cylinder/acyclic-closure work and Levine's homology-cylinder/tree-level work without precise bibliographic entries or resolving theorems.
\end{itemize}
\end{document}
